% RESEARCH_DRAFT_SHA256: 96d6b1285938b03c35f0a323689a87fe54fb3d32b8c5b92a5a640f116664d9a9
% Layout and source-reference wrappers expanded for the web renderer.
\documentclass[11pt]{article}
\usepackage[margin=25mm]{geometry}
\usepackage{amsthm}
\usepackage{fontspec,unicode-math}
\setmainfont{Latin Modern Roman}
\setmathfont{Latin Modern Math}
\usepackage{amsmath,amssymb,mathtools}
\usepackage{xurl,hyperref}
\hypersetup{colorlinks=true,urlcolor=blue,linkcolor=blue}
\setcounter{secnumdepth}{0}
\setlength{\parindent}{0pt}
\setlength{\parskip}{5pt}
\setlength{\emergencystretch}{3em}
\newcommand{\sref}[2]{\hyperlink{source:#1:#2}{[S#2]}}
\newcommand{\cataloguescope}[1]{\paragraph{Source question.}#1}
\newcommand{\floor}[1]{\left\lfloor#1\right\rfloor}
\newtheorem{proposition}{Proposition}
\newtheorem{lemma}[proposition]{Lemma}
\begin{document}
\paragraph{Research manuscript.} OpenAI \texttt{gpt-6-astra}, reasoning \texttt{ultra}; 11 October 2026.
\paragraph{Preserved source body.} The complete catalogue mathematical text and sources follow with layout and source links expanded, including source qualifications and its historical preparation notes. New research begins at the section headed Research attempt.

\section{Root obstruction and outgroup-oriented algebraic clade certificates}\label{statement}

\subsection{Definitions and mathematical statement}

\paragraph{Shared notation.}Standard number systems are denoted by $\mathbb N,\mathbb Z,\mathbb Q,\mathbb R,\mathbb C$; index ranges, computational models and specialized notation follow the cited problem source and the definitions below.


A phylogenetic model assigns transition parameters to edges of a tree and maps them to joint probabilities at its observed leaves. A phylogenetic invariant is a polynomial vanishing on all probability vectors in that model. A clade is a set of descendants of one node in a rooted phylogenetic tree. Identifying a clade from leaf data requires a rooted-tree convention or an outgroup. Algebraic invariant techniques test polynomial relations on model probability distributions.
\paragraph{Statement recorded in the supplied source.}
Question 36. Can we develop algebraic invariant techniques for identifying clades? \href{https://aimath.org/WWN/compalgstat/compalgstat.pdf}{[S2]}
\paragraph{Scope qualification.}
The supplied extraction includes adjacent workshop exposition; the numbered target passage is retained verbatim and that surrounding exposition is omitted. The original question does not stipulate a mutation model or rooting convention; unrooted splits and rooted clades are distinct tasks.
\paragraph{Source status.}
The supplied catalogue labels this record partially solved. Its source-status metadata was checked against the supplied snapshot on 2026-10-11. This import does not independently establish current open status or mathematical correctness.
\subsection{Short English statement}
Can algebraic invariants help identify groups of species that descend from one ancestral node?
\subsection{Sources}
\begin{itemize}
\item[S1] Ulam AI and the UnsolvedMath Contributors, UnsolvedMath: A Curated Collection of Open Mathematics Problems, catalogue entry. Catalogue curation and metadata: CC BY 4.0. The supplied source question is retained; definitions and the English summary are editorial additions.
\url{https://huggingface.co/datasets/ulamai/UnsolvedMath}.
\item[S2] AIM Workshop Problem Lists, original problem source linked by the catalogue
\url{https://aimath.org/WWN/compalgstat/compalgstat.pdf}.
\end{itemize}
\paragraph{Preparation and review.}
Original question attributed to the cited source and UnsolvedMath contributors. Codex (GPT-6; exact serving revision not exposed) assisted with definitions, English summaries, source inspection and formatting. No new human mathematical authorship or independent proof verification is asserted. Catalogue identity checked against the complete supplied registry and existing repository statements; known flagged repetitions and obvious identical conjectures excluded. Full mathematical-equivalence/literature review has not been performed. Statement preparation and final review followed contribution prompts C and D; automated checks verify structure and identity, not open status.
\label{end:statement}


\section{Research attempt}
\subsection{Literature and scope}
Question 36 asks about identifying clades without specifying root information or a mutation model. The catalogue's historical discussion invokes reversible models. I tested the suggested escape of allowing general nonreversible transition matrices. With unrestricted positive edge matrices, rerooting still prevents identifying rooted clades from the leaf distribution.
\begin{proposition}
Every strictly positive general Markov distribution on a finite tree has a representation rooted at any other vertex with positive row-stochastic edge matrices and the same full joint distribution.
\end{proposition}
\begin{proof}
Move the root across one edge from $u$ to $v$. Let $\pi_a>0$ be the root law at $u$ and $M_{ab}>0$ its transition to $v$. Set
\[
\pi'_b=\sum_a\pi_aM_{ab},\qquad K_{ba}=\frac{\pi_aM_{ab}}{\pi'_b}.
\]
Then $\pi'_b>0$ and $\sum_aK_{ba}=1$. In the factorized full joint probability, replace the factor $\pi_aM_{ab}$ by $\pi'_bK_{ba}$, which is identical. All other edges retain their orientations when the root crosses this single edge, and all their factors are unchanged. Thus the full joint law, and hence every leaf marginal, is identical. Repeat along the path to the desired root.
\end{proof}
\subsection{Failed inference and a concrete instance}
On a quartet with internal vertices $u,v$, leaves $1,2$ attached to $u$ and $3,4$ to $v$, root $u$ makes $\{3,4\}$ a descendant clade whereas root $v$ makes $\{1,2\}$ one. These differing rooted interpretations have the same leaf law after the above reparameterization. For $\pi=(3/5,2/5)$ and $M=\left(\begin{smallmatrix}4/5&1/5\\1/3&2/3\end{smallmatrix}\right)$, the new root law is $(46/75,29/75)$ and reverse matrix is $\left(\begin{smallmatrix}18/23&5/23\\9/29&20/29\end{smallmatrix}\right)$. Direct multiplication verifies every joint edge probability.
\subsection{Verification and remaining gaps}
The rational example checks positivity, row sums and the identity $\pi_aM_{ab}=\pi'_bK_{ba}$. This is a standard probabilistic reparameterization, not a new identifiability theorem. It shows that nonreversibility by itself is insufficient when matrices are unrestricted. Shared edge parameters, clocks, outgroups or other restrictions may obstruct rerooting; no claim about those models is made. Rooted-clade identification under suitable specified assumptions remains to be developed.

\section{References and source attribution}
The original question is credited to the AIM workshop list prepared by Seth Sullivant (2004), \url{https://aimath.org/WWN/compalgstat/compalgstat.pdf}. The source question and catalogue attribution are reproduced above. This list was accessed on 11 October 2026; it is historical evidence of the question, not a certification of present open status.


\paragraph{Integration review.} Codex (GPT-6; exact serving revision not exposed) reviewed the written arguments and expanded layout and source-reference wrappers. Repository metadata and the production web renderer were checked using Python 3.13 and JavaScript with MathJax 3.2.2. This model-assisted review is not independent human proof verification; no PDF compilation or formal proof checking is asserted.

\section{Completion Estimate}
\noindent\textbf{COMPLETION: 0\%}
No substantiated new progress toward the entire original target is claimed. The worked results and failed route are delimited above; a known result, illustrative model, or auxiliary obstruction is not counted as original-target completion.

\section{Final status and attribution}
\textbf{Attempt status: unresolved.} This describes this research attempt, not a new ruling on the source catalogue's status or on external literature claims.
The mathematics and this self-audit were generated by OpenAI \texttt{gpt-6-astra}, reasoning \texttt{ultra}, on 11 October 2026. Codex (GPT-6; exact serving revision not exposed) orchestrated the batch and contributed formatting/repository review. No human mathematical authorship, independent human proof checking, or formal verification is claimed. The historical preparation paragraphs in the reproduced catalogue body concern that body only.

Contribution prompt A was applied to the complete definition and supplied source record, followed by the adversarial checks of prompt B. Local reproduction files are \texttt{batch005/checks.py}, \texttt{computations.json}, and the per-ID audit JSON. They were not public links at preparation time. Numerical checks are finite and do not prove unrestricted results. No LaTeX engine was found locally; compilation was not performed by this research agent.

\end{document}
